% CVPR 2026 Paper Template; see https://github.com/cvpr-org/author-kit

\documentclass[10pt,twocolumn,letterpaper]{article}

%%%%%%%%% PAPER TYPE  - PLEASE UPDATE FOR FINAL VERSION
\usepackage{cvpr}              % To produce the CAMERA-READY version
% \usepackage[review]{cvpr}      % To produce the REVIEW version
% \usepackage[pagenumbers]{cvpr} % To force page numbers, e.g. for an arXiv version

% Import additional packages in the preamble file, before hyperref
\input{preamble}

% It is strongly recommended to use hyperref, especially for the review version.
% hyperref with option pagebackref eases the reviewers' job.
% Please disable hyperref *only* if you encounter grave issues, 
% e.g. with the file validation for the camera-ready version.
%
% If you comment hyperref and then uncomment it, you should delete *.aux before re-running LaTeX.
% (Or just hit 'q' on the first LaTeX run, let it finish, and you should be clear).
\definecolor{cvprblue}{rgb}{0.21,0.49,0.74}
\usepackage[pagebackref,breaklinks,colorlinks,allcolors=cvprblue]{hyperref}

\usepackage{verbatim}
\usepackage{tikz}
\usetikzlibrary{positioning}

%%%%%%%%% PAPER ID  - PLEASE UPDATE
\def\paperID{Project 07} % *** Enter the Paper ID here
\def\confName{CVPR}
\def\confYear{2026}

%%%%%%%%% TITLE - PLEASE UPDATE
\title{Compact Feed-Forward 3D Reconstruction through Distillation and Test-Time Optimization}

%%%%%%%%% AUTHORS - PLEASE UPDATE
\author{
Antonin Granados\\
{\tt\small antonin.granados@telecom-paris.fr}
\and
Finn Ferchau\\
{\tt\small finn.ferchau@tum.de}
}

\begin{document}
\maketitle
\input{fig/timeline}
\begin{abstract}
Feed-forward multi-view reconstruction models such as VGGT and DéjàView predict scene geometry and camera parameters in a single pass, but achieving high reconstruction quality can require substantial computational resources. We investigate whether test-time optimization (TTO) can recover the performance lost when distilling large reconstruction models into smaller students. We will train progressively compressed students from teacher predictions, refine their outputs using multi-view consistency, and compare them with the original models at matched computational budgets. The evaluation will characterize the trade-off between model size, reconstruction quality, and optimization cost, identifying which initialization errors TTO can correct and the smallest student that approaches teacher-level performance.

\begin{comment}
\begin{itemize}
    \item Feed-forward multi-view transformers (VGGT, DéjàView) predict cameras, depth and point maps in a single pass, but state-of-the-art models have grown to $\sim$1B parameters
    \item Two observations: compact models can be competitive (DéjàView, 117M), and optimization on top of feed-forward predictions still gives large gains (VGGT + BA)
    \item Question: how small can a feed-forward model become if test-time optimization (TTO) refines its output? Does high-quality reconstruction need an accurate initial prediction?
    \item Approach: distill VGGT/DéjàView into progressively smaller students, refine their predictions with TTO, and compare against the teachers
    \item Expected result: quality-vs-compute Pareto curves and a characterization of which initialization errors TTO can and cannot correct
\end{itemize}
\end{comment}

\end{abstract}    
\section{Introduction}
\label{sec:intro}

\subsection{Motivation}
Classical structure-from-motion (SfM) pipelines, such as COLMAP~\cite{sfm}, reconstruct 3D scenes through feature matching, camera pose estimation, triangulation and bundle adjustment (BA). While effective, these pipelines can be slow and fail when there is little texture, limited overlap or sparse view collections.

Feed-forward multi-view transformers, such as DUSt3R~\cite{dust3r} and VGGT~\cite{vggt}, offer an alternative by predicting camera parameters and 3D point maps directly from images. These models can provide rebust initial reconstruction. However, they are accompanied by a substantial model size: VGGT for instance contains around 1.26B parameters.

\begin{itemize}
    \item Classical SfM (COLMAP \cite{sfm}): matching, pose estimation, triangulation and bundle adjustment (BA); accurate when it converges, but slow and brittle with little texture, little overlap or few views
    \item Feed-forward multi-view transformers (DUSt3R \cite{dust3r}, VGGT~ \cite{vggt}) regress cameras, depth and point maps in one pass and stay robust where SfM fails
    \item Progress came mostly from scaling: VGGT has 1.26B parameters
    \item At similar size, the fully trained DéjàView roughly matches $\pi^3$, while the cheaply distilled OTT3R student falls short \cite{dejaview,ott3r}. At small training budgets the limit is training compute. TTO spends compute per scene instead
    \item Gap: model compression and TTO have so far been studied separately; how they trade off at a fixed compute budget is open \todo{double-check with a literature search}
\end{itemize}

\subsection{Research Questions}
\begin{itemize}
    \item \textbf{RQ1:} How small can a distilled student become while still providing an initialization from which TTO reaches teacher-level quality?
    \item \textbf{RQ2:} At matched total compute, is a compact model with TTO competitive with a large feed-forward model?
\end{itemize}

\subsection{Contributions (planned)}
\begin{itemize}
    \item A family of compact students distilled from VGGT/DéjàView 
    \item A compute-matched study of student size $\times$ TTO budget with Pareto curves against the teachers
    \item one methodological contribution. Hard to plan at this point in time.
\end{itemize}
\section{Related Works}
\label{sec:related}

\subsection{Multi-View 3D Reconstruction}
\begin{itemize}
    \item Classical SfM, \eg COLMAP \cite{colmap}: feature matching, incremental registration, triangulation and BA; accurate but slow and brittle
    \item DUSt3R \cite{dust3r}, MASt3R \cite{mast3r}: regress point maps for image pairs, then fuse all pairs with a global alignment optimization
    \item VGGT \cite{vggt}: 1.2B-parameter transformer with alternating frame/global attention over DINOv2 tokens; predicts cameras, depth, point maps and tracks for up to hundreds of views in one pass
    \item $\pi^3$ \cite{pi3}, Depth Anything 3 \cite{da3}: further large feed-forward models (959M; 356M--1.2B)
    \item DéjàView \cite{dejaview}: a single weight-shared transformer block applied $K$ times; 117M parameters, matches or beats $\sim$1B-parameter models, $K$ is a compute knob at inference
\end{itemize}

\subsection{Reconstruction via Distillation}
\begin{itemize}
    \item Knowledge distillation: a student is trained to reproduce a teacher's outputs \cite{hintondistillation}
    \item OTT3R \cite{ott3r}: distills $\pi^3$ (959M) into a 102M student (9.4$\times$ smaller, up to 7$\times$ faster) using 1.6\% of VGGT's training compute; also a pseudo-labeling pipeline as a fast alternative to COLMAP. Closest prior work, but purely feed-forward \todo{confirm no TTO in full paper}
    \item Multi-View 3D Reconstruction using Knowledge Distillation \cite{reconstructiondistillation}: distills DUSt3R into scene-specific CNN/ViT students on 12Scenes; does not generalize to new scenes
    \item A survey on knowledge distillation: Recent advancements \cite{surveydistillation} \todo{(mostly for us)}
\end{itemize}

\subsection{Test-Time Optimization}
\begin{itemize}
    \item VGGT + BA \cite{vggt}: BA initialized from VGGT's predicted points and tracks, no separate triangulation needed
    \item Test3R \cite{test3r}: test-time training for DUSt3R-style pairwise models; maximizes the consistency of the reconstructions from $(I_1, I_2)$ and $(I_1, I_3)$ relative to the shared image
    \item Learning 3D Reconstruction with Priors in Test Time \cite{reconstructionpriors}: optimizes a pretrained multi-view transformer at inference with self-supervised cross-view photometric/geometric consistency, plus optional priors (poses, intrinsics, depth) as penalty terms; code is public
    \item VGG-$T^3$ \cite{vggt3}: uses test-time training as an architectural component, compressing the growing key-value representation into a fixed-size MLP so that cost scales linearly with the number of views; targets efficiency, not output refinement
    \item Positioning: these works refine large models; we study TTO as a substitute for model capacity
\end{itemize}

\section{Approach}
\label{sec:approach}

\subsection{Problem Formulation}
\todo{notations taken from VGGT \cite{vggt}}

Sequence $(I_i)_{i=1}^N$ of $N$ RGB images $I_i \in \mathbb{R}^{3 \times H \times W}$ capturing the same 3D scene. The goal is to reconstruct the scene by finding each camera's (i.e., for each image) intrinsics and extrinsics $\mathbf{g}_i \in \mathbb{R}^9$, as well as a point map $P_i \in \mathbb{R}^{3 \times H \times W}$ (which can be computed from the camera parameters and a depth map $D_i \in \mathbb{R}^{H \times W}$).

\begin{itemize}
    \item Feed-forward model $f_\theta: (I_i)_{i=1}^N \mapsto (\mathbf{g}_i, D_i, P_i)_{i=1}^N$; teacher $f_{\theta_T}$ (VGGT or DéjàView), student $f_{\theta_S}$ with $|\theta_S| \ll |\theta_T|$
    \item TTO $\mathcal{O}_T$ refines the student's prediction for $T$ iterations using only the input images: $(\hat{\mathbf{g}}_i, \hat{D}_i)_{i=1}^N = \mathcal{O}_T\big(f_{\theta_S}(I), I\big)$
    \item Total cost $C = C_{\text{fwd}}(\theta_S) + T \cdot C_{\text{iter}}$ (FLOPs or wall-clock time); for a budget $C$, which (student size, $T$) minimizes the reconstruction error, and how does it compare to $f_{\theta_T}$ at the same $C$?
    \item If DéjàView is the teacher, the model predicts depth and ray maps; $\mathbf{g}_i$ and $P_i$ are derived from the rays \cite{dejaview}
\end{itemize}

\subsection{Proposed Method}
\begin{itemize}
    \item \textbf{Teacher} \todo{decide}: VGGT (1.26B, public code, predicts tracks usable for BA) and/or DéjàView (117M, stronger, already compact)
    \item \textbf{Students:} progressively smaller variants of the teacher architecture
    \begin{itemize}
        \item smaller DINOv2 encoder (ViT-B, 86M $\rightarrow$ ViT-S, 21M), initialized from pretrained weights
        \item fewer or narrower transformer blocks (VGGT) or fewer loop steps $K$ (DéjàView)
        \item \todo{size grid, \eg 4 students between $\sim$20M and $\sim$300M}
    \end{itemize}
    \item \textbf{Distillation:}
    \begin{itemize}
        \item output level: student regresses the teacher's depth, rays/cameras and point maps with the teacher's training losses, weighted by the teacher's confidence
        \item optional feature level: match intermediate tokens (needs a projection when widths differ)
        \item optionally mix in ground-truth losses where available (ablation)
    \end{itemize}
    \item \textbf{Test-time optimization} (per scene, no ground truth, initialized from the student):
    \begin{itemize}
        \item variables: camera parameters $\mathbf{g}_i$ and per-view depth corrections (\eg scale/shift + low-resolution residual) \todo{or network weights as in \cite{test3r,reconstructionpriors}}
        \item objectives: reprojection error on correspondences (BA; tracks from an off-the-shelf matcher or from the predicted point maps), photometric or feature-metric consistency between views warped with the predicted depth and poses, multi-view depth consistency, and a regularizer towards the initialization
        \item number of iterations $T$ is the second compute knob next to student size
    \end{itemize}
    \item \textbf{Candidate contributions} \todo{pick 1--2}:
    \begin{itemize}
        \item optimization-aware distillation: apply the loss after a few unrolled, differentiable TTO steps so that the student spends its capacity on errors TTO cannot fix
        \item learned vs.\ hand-designed iteration: DéjàView loop steps $K$ against TTO iterations at equal compute
        \item confidence-guided TTO: use predicted confidence to decide where and how long to optimize
    \end{itemize}
\end{itemize}

\subsection{Training (dataset ?)}
\begin{itemize}
    \item Teacher pseudo-labels remove the need for 3D ground truth $\rightarrow$ any multi-view images or videos can be used
    \item Candidate data: subsets of the teachers' training sets that also have ground truth, \eg CO3Dv2, DL3DV, BlendedMVS, MegaDepth, ScanNet++ (train split), Hypersim, Virtual KITTI 2 \cite{vggt,dejaview}; OTT3R's pseudo-labeling pipeline as an alternative \cite{ott3r}
    \item Teacher outputs precomputed offline or produced online: depth + rays + confidence at $\sim$500\,px $\approx$ 3\,MB per image in fp16, \ie $\sim$300\,GB per 100k images \todo{decide}
    \item Reduced-scale training: 2--24 views per sample, $\sim$500\,px longest side \cite{vggt}; encoder frozen or trained at a lower learning rate (DéjàView uses 0.1$\times$ on DINOv2 \cite{dejaview})
    \item Evaluation scenes stay out of the distillation data (scene-level split for ScanNet++)
    \item \todo{ask supervisors: available GPUs and storage}
\end{itemize}

\section{Evaluation}
\label{sec:eval}

\subsection{Baselines}

\subsection{Metrics}
\section{Plan}
\label{sec:plan}

The project timeline is shown in Fig.~\ref{fig:plan}.

\subsection{Milestones}

\begin{itemize}
    \item \textbf{M1 - Baselines:} run VGGT and DéjàView on the benchmarks; reproduce reported numbers
    \item \textbf{M2 - Distillation:} training pipeline and pseudo-labels; train students of increasing compression; evaluate their initial predictions
    \item \textbf{M3 - TTO (starts in parallel with M1):}
    \item \textbf{M4 - Comparison and writing:} matched-compute comparison, ablations, chosen methodological contribution, paper
\end{itemize}

\subsection{Expected Outcome}

The project will produce Pareto curves of reconstruction quality versus total compute, comparing distilled students with TTO against the original teachers. We will identify the smallest student for which TTO approaches teacher-level quality (Q1), and release compact student checkpoints and evaluation code. Even if TTO cannot compensate for reduced model capacity, the results will help quantify the initialization quality required for effective refinement.

\begin{comment}
\begin{itemize}
    \item Pareto curves of reconstruction quality vs.\ total compute for students + TTO against the teachers
    \item The smallest student from which TTO reaches teacher quality (RQ1)
    \item Compact student checkpoints and code
    \item A negative result (TTO cannot replace capacity) would still be informative: it quantifies how accurate the initialization has to be
\end{itemize}
\end{comment}

{
    \small
    \bibliographystyle{ieeenat_fullname}
    \bibliography{main}
}

\end{document}
